\documentclass[11pt,reqno]{amsart}

\usepackage[T1]{fontenc}
\usepackage{lmodern}
\usepackage{amsmath,amssymb,amsthm,mathtools}
\usepackage{mathrsfs}
\usepackage[colorlinks=true,linkcolor=blue,citecolor=blue,urlcolor=blue]{hyperref}
\usepackage[nameinlink,noabbrev]{cleveref}

\newtheorem{theorem}{Theorem}[section]
\newtheorem{proposition}[theorem]{Proposition}
\newtheorem{lemma}[theorem]{Lemma}
\newtheorem{corollary}[theorem]{Corollary}
\theoremstyle{definition}
\newtheorem{definition}[theorem]{Definition}
\newtheorem{remark}[theorem]{Remark}

\newcommand{\Q}{\mathbb Q}
\newcommand{\Z}{\mathbb Z}
\newcommand{\R}{\mathbb R}
\newcommand{\C}{\mathbb C}
\newcommand{\Hh}{\mathfrak H}
\newcommand{\OO}{\mathcal O}
\newcommand{\Cl}{\operatorname{Cl}}
\newcommand{\trd}{\operatorname{trd}}
\newcommand{\nrd}{\operatorname{nrd}}
\newcommand{\rad}{\operatorname{rad}}
\newcommand{\vol}{\operatorname{vol}}
\newcommand{\Fitt}{\operatorname{Fitt}}
\newcommand{\dd}{\mathbf d}
\newcommand{\eps}{\varepsilon}
\newcommand{\Kc}{\mathcal K}

\title[The $5/3$ valuation bound for $abc$ triples]{An Aggregate CM-Pair Estimate and a $5/3$ Bound for Valuation Products of $abc$ Triples}
\author{Ruiyi Zhang}
\thanks{Independent Researcher, Columbia, Maryland, USA.  Email:
\href{mailto:ryzfsu@gmail.com}{ryzfsu@gmail.com}.}
\date{July 2026}
\subjclass[2020]{Primary 11G18; Secondary 11G05, 11D75, 11R29}
\keywords{abc triples, Shimura curves, CM points, quaternion orders, valuations}

\begin{document}

\begin{abstract}

Let $a+b=c$ be pairwise coprime positive integers. Then, we introduce $R=\rad(abc)$ and $P(abc)=\prod_{p\mid abc}v_p(abc)$. We present here the proof of \[
 P(abc),\ \dd(abc)\ll_{\eps}R^{5/3+\eps},
\]
where $\dd$ is taken to be the divisor function. The first step in the deduction gives an aggregate fixed-radius estimate for CM points on an indefinite Shimura curve. In the case where the CM discriminants lie in $[\Delta,2\Delta]$, the number of nearby ordered pairs is \[
 \#\{(P,Q)\in\mathscr P_\Delta^2:d(P,Q)\leq r_0\}
 \ll |\mathscr P_\Delta|+
 \Delta\left(1+\frac{\Delta^2}{N}\right)N^\eps,
\]
which holds uniformly in all polynomial ranges $\Delta\leq N^A$. This is established by an elementary local orbit count in the Bruhat--Tits tree, using the characteristic-matrix framework of Bruinier, Li and M\"oller. The spatial Cauchy--Schwarz then yields a bound for the first absolute moment of the CM-packet. The ordinary AM--GM now applies directly to the logarithmic average in Pasten's Arakelov estimate to complete the argument, and neither a fourth moment nor a pointwise subconvexity estimate is required.

\end{abstract}
\maketitle
\pagestyle{plain}


\newpage

\section{Introduction}

Let $a,b,c$ be pairwise coprime positive integers satisfying $a+b=c$. Here \[
 R=\rad(abc),\qquad
 P(abc)=\prod_{p\mid abc}v_p(abc),\qquad
 \dd(n)=\prod_{p\mid n}\bigl(v_p(n)+1\bigr).
\] Using a comparison between classical and quaternionic modular degrees,
Pasten was able to derive the unconditional bounds $P(abc),\dd(abc)\ll_\eps R^{8/3+\eps}$
by evaluating an integral differential at CM points \cite[Theorem 1.11 and
Theorem 16.8]{Pasten}. 
Our main result is a reduction of the exponent.

Rickards has established precise formulas regarding the intersections of
closed geodesics associated with pairs of real quadratic embeddings. These formulas include $x$-linking and local Eichler-order information
\cite[Theorems 10 and 63]{Rickards}. More recently, Bruinier--Li--M\"oller considered all signatures of a pair of
quadratic vectors. More specifically, their formula gives the number of pairs of Heegner points
at a given distance, and this number is a product of local factors
\cite[Theorem 1.1(iv) and Corollary 5.9]{BruinierLiMoller}. We use Bruinier--Li--M\"oller's characteristic-matrix and distance
formalism, but prove the coarse orbit bound needed here directly from the
local structure of quaternion orders and the Bruhat--Tits tree
\cite[Equation (23.2.4), Theorem 13.3.11, and Sections
23.5.11--23.5.16]{Voight}. The aggregate estimate provides us with the divisor-sum bound we need for
Pasten's argument.

In the context of the theory of special cycles, Heegner points, and
quaternionic Shimura varieties, the natural objects in that framework are CM
points
\cite{KudlaRapoportYang,ZhangCMEquidistribution,ZhangHeegner,YuanZhangColmez}. 
Complete or structured packets satisfy the equidistribution, variance, and
small-scale bounds
\cite{MichelHeegner,HarcosMichel,Khayutin,LiuMasriYoung,AlsinaChatzakos, HumphriesRadziwill,WieserYang}. None of these works states our fixed-radius bound for arbitrary selected
submultisets. 
Our approach is to substitute the classical sup-norm techniques
\cite{BlomerHolowinsky,DasSengupta} with a local $L^1$ inequality. Later abc work treats different functionals or additional size hypotheses
\cite{PastenSepulveda}.

\begin{theorem}\label{thm:main}
For every $\eps>0$ and for every pairwise coprime triple of positive integers
$a+b=c$, we have
\[
 P(abc),\ \dd(abc)\ll_\eps \rad(abc)^{5/3+\eps}.
\]
\end{theorem}

Geometrically speaking, let $R_N$ stand for an Eichler order of reduced
discriminant $N$ in an indefinite quaternion division algebra, putting
$X_N=R_N^1\backslash\Hh$. We assume that the odd part of $N$ is square-free, and that $v_2(N)\leq8$. For odd fundamental discriminants, $-d$ being coprime to $N$, and
$d\in[\Delta,2\Delta]$, choose $\mathscr P_\Delta$ as any multiset of the
related CM centers where the multiplicity of each one does not exceed a
fixed $B_0\geq1$.

\begin{theorem}[Aggregate CM-pair estimate]\label{thm:pair-intro}
Fix $A,r_0>0$ and $B_0\geq1$.  If $\Delta\leq N^A$ and every multiplicity
in $\mathscr P_\Delta$ is at most $B_0$, then
\[
 \#\{(P,Q)\in\mathscr P_\Delta^2:d(P,Q)\leq r_0\}
 \ll_{A,r_0,B_0,\eps}
 |\mathscr P_\Delta|+
 \Delta\left(1+\frac{\Delta^2}{N}\right)N^\eps.
\] The estimate is uniform both in the chosen discriminant submultiset and
in the admissible factorization $N=DM$.
\end{theorem}

For pairs lying on the off-diagonal, we will set $-d,-e$ to be the discriminants and $b$ to be the integral pairing parameter. The local conditions state that the odd part of $N$ has to divide $b^2-de$. A direct local orbit count gives an upper bound of $N^\eps$ for each global orbit count. After we have fixed the radius we will be down to $O_{r_0}(\Delta(1+\Delta^2/N))$ choices for $b$ and the quotient. These choices fix $de$, and $de$ can have at most $N^\eps$ admissible factorizations. Further, the thinning of the packets considered here will only eliminate pairs.


In the case of $\mathscr P$, that is, the corresponding full Galois CM
multiset from Pasten's construction, we use $x=N^{2/3+\delta}$ and
$\vartheta=\delta/5$. Then, using \cref{thm:pair-intro}, the local holomorphic $L^1$ inequality,
and the spatial form of the Cauchy--Schwarz inequality, we obtain \[
 \frac1H\sum_d\frac1{2h_d}
 \sum_{P\in\mathscr C_d}|F_N(P)|
 \ll N^{\tau+\eps},
 \qquad
 \tau=\left(\frac23+\delta\right)\frac{6\delta}{5}.
\]
Here $\|F_N\|_{L^2(\mu_{\mathrm{prob}})}=N^{o(1)}$. Applying zero-counting and Markov's inequality at this point, we obtain a full
packet that is free of zeros, on which the first absolute moment is
$\ll N^{\tau+\eps}$. We write $u_\sigma=|f(\tau_\sigma)|\operatorname{Im}\tau_\sigma>0$, and for these we need only AM--GM:
\[
 \frac1{2h_d}\sum_\sigma\log u_\sigma
 \leq \log\left(\frac1{2h_d}\sum_\sigma u_\sigma\right).
\]

Using these first moment and AM--GM bounds along with Pasten's Proposition
14.2, height formula, and modular-degree comparison gives the cell exponent
$q_\delta=5/3+13\delta/5+12\delta^2/5$. In place of Pasten's norm exponent value of $5/6$, we can put the value
$1/3+o(1)$ that our aggregate first-moment estimate gives us. With that replacement, $1+2(5/6)=8/3$ changes to $1+2(1/3)=5/3$. Finally, the cyclic pairs and triples will distribute the level factors that
are left, and we let $\delta\to0$ to obtain the proof of
\cref{thm:main}.

\section{Quaternionic CM packets}\label{sec:setup}

Let $B$ represent an indefinite quaternion division algebra over $\Q$, and let $R_N$ denote an Eichler order of reduced discriminant $N=DM$, where $D$ is the discriminant of $B$ and $M$ the Eichler level. As in Pasten, a factorization $N=DM$ is called admissible when $(D,M)=1$ and $D$ is the product of an even number of distinct primes; in the compact case considered below, $D>1$. We assume that the odd part of $N$ is square-free and that $v_2(N)\leq8$. Write
\[
 \Gamma_N=R_N^1/\{\pm1\},\qquad X_N=\Gamma_N\backslash\Hh.
\]
In this notation we write $d\mu_{\mathrm{phys}}=dx\,dy/y^2$ for the
hyperbolic measure in the physical sense, we let $\mu_{\mathrm{prob}}$ be
the probability measure, and we set $V_N=\vol(X_N)$. Using the exact volume formula from \cite[Equation (2.7)]{BruinierLiMoller},
we get, for every $\eps>0$, \[ N^{1-\eps}\ll_\eps V_N\ll_\eps N^{1+\eps}.
\]

As for the Gross lattice, we can write it out in the following form
\[
 L_N=\{2a-\trd(a):a\in R_N\}\subset B^0.
\] In what follows we adopt the source convention $Q=-\nrd$ on the
trace-zero space, with its full polarization \[ B_Q(x,y)=Q(x+y)-Q(x)-Q(y) \]
on the same space, so that $B_Q(x,y)=\trd(xy)$.
We note that the primitive vectors representing imaginary quadratic
discriminant $-d$ satisfy $Q(x)=-d<0$. 
Taking lifts and signs for the associated negative lines, we read off the
relation \[ \frac{|B_Q(x,y)|}{2\sqrt{Q(x)Q(y)}}=\cosh d(x,y).
\] In fact, if we divide $x$ and $y$ by the square roots of $-Q(x)$ and
$-Q(y)$, we recognize the standard distance equation for the negative sheet
of the hyperboloid model, the $2$ being the contribution of the full
polarization.

\begin{lemma}[Lattice and characteristic-matrix normalization]
\label{lem:normalization} The lattice $L_N$ is exactly the index-four
lattice $L$ that we find in \cite[Equation (2.9)]{BruinierLiMoller}. An optimal embedding of the imaginary quadratic
maximal order of odd fundamental discriminant $-d$ determines, up to sign, a
primitive $x\in L_N$ with $Q(x)=-d$, and the converse holds as well. For a
pair $x,y$, the characteristic matrix in the notation of
Bruinier--Li--M\"oller is
\[
 T=\frac12\bigl(B_Q(x_i,x_j)\bigr)_{i,j}
  =\begin{pmatrix}-d&b\\ b&-e\end{pmatrix},
 \qquad b=\frac12B_Q(x,y).
\] \end{lemma} \begin{proof} If $a\in R_N$, we can check that either
$2a-\trd(a)$ lies in $2(R_N\cap B^0)$ or $(1+2a-\trd(a))/2\in R_N$, 
according to the parity of $\trd(a)$. In both cases the converse holds as well, and we obtain the equality with
the lattice in \cite[Equation (2.9)]{BruinierLiMoller}. Let us assume now that $x\in L_N$ and $Q(x)=-d$, with $-d$ being an odd
fundamental discriminant. Then $x$ cannot belong to $2(R_N\cap B^0)$, therefore $\alpha=(1+x)/2\in R_N$. Since \[ \alpha^2-\alpha+\frac{d+1}{4}=0, \] $\alpha$ defines an
embedding of the maximal order $\Z[(1+\sqrt{-d})/2]$ into $R_N$, and this
embedding is optimal since the source order is maximal.
Conversely, an optimal embedding $\varphi$ sends $\sqrt{-d}$ to the
vector $x=\varphi(\sqrt{-d})\in L_N$ with $Q(x)=x^2=-d$. This vector is primitive.
Indeed, $Q(L_N)\subset\Z$ because $Q(2a-\trd(a))=\trd(a)^2-4\nrd(a)$ for
$a\in R_N$, and hence a factorization $x=kx_0$ in $L_N$, with $k>1$, would
force $k^2\mid d$, contrary to the square-freeness of $d$. This proves the
stated correspondence, and it also proves the normalization in the
characteristic matrix $\frac12(B_Q(x_i,x_j))$. In addition, let us state the values $x=2\alpha-1$ and $y=2\beta-1$, where
$\alpha,\beta\in R_N$ both have the reduced trace value of one.
Then
\[
 B_Q(x,y)=4\trd(\alpha\beta)-2,
 \qquad b=2\trd(\alpha\beta)-1\in\Z,
\]
so $b$ is odd. This direct calculation fixes the normalization that we will use below, and
we can check that its off-diagonal entry agrees with \cite[Equation
(2.14)]{BruinierLiMoller}.
\end{proof} In what follows, we will take CM discriminants to be coprime to
$N$ and to lie in $[\Delta,2\Delta]$.

\begin{remark}\label{rem:galois-packet} In the following, we will write
$h_d$ for the class number of $\Q(\sqrt{-d})$, while $\mathscr C_d$ will
denote the complete $2h_d$-element Galois multiset of Pasten, and not a
selected orientation packet. Moreover, we can use the fact that the CM point of Pasten has a residue
field equal to the Hilbert class field \cite[Section 13.1]{Pasten}, which
ensures that these $2h_d$ conjugates are distinct quotient points. Each of the elements is a CM center of discriminant $-d$, and hence we can
see that it has the primitive-vector realization in
\cref{lem:normalization}. Since the pair estimate is uniform for submultisets, we do not need to
identify this Galois multiset with one prescribed system of local
orientations. As for standard optimal-embedding counts with a fixed local class, we can
find such counts stated in \cite[Theorem 2.5]{PadurariuSaia}.
\end{remark} \begin{remark}\label{rem:pair-formula} For nonsingular $T$, we
can recall that Bruinier--Li--M\"oller identify $R_T$ with the diagonal
norm-one orbit set of ordered pairs in $L_N^2$ that have characteristic
matrix $T$. See their
definition of $R_T$ together with \cite[Equation
(2.9)]{BruinierLiMoller}. The paragraph preceding their Theorem~5.7 explicitly states that the lattice
in use for their subsequent local computations is reset to the localization
of this same index-four lattice. Their count allows for the inclusion of imprimitive pairs while the CM
centers are derived from primitive vectors, so we can see that the latter
classes form a subset. The bound for the entire orbit set below is obtained
from a direct local coset count and does not use an exact local-density
formula.
Indeed, a CM center determines its negative rational line, on which a
primitive vector of norm $-d$ is unique up to sign, since any two such
vectors differ by a rational scalar with square equal to one. In moving to the quotient, at most an elliptic
stabilizer can be added, and by the integral-trace argument in the proof
of \cref{lem:local-l1}, the projective order of that stabilizer is bounded by $3$. Thus the fibers of the map from primitive-vector
classes to CM centers are uniformly bounded, and the bounded multiplicities
in $\mathscr P_\Delta$ will be absorbed in the implied constant.
\end{remark}

\section{The aggregate CM-pair estimate}\label{sec:pair}

We prove \cref{thm:pair-intro}.

\begin{lemma}[Primitive local-factor bound]\label{lem:local-factors}
Let
\[
 T=\begin{pmatrix}-d&b\\ b&-e\end{pmatrix},
 \qquad N^\flat=N/2^{v_2(N)},
\]
where $-d,-e$ are odd fundamental discriminants, $(de,N)=1$, $b$ is
odd, and $T$ is nonsingular. Assume that the odd part of $N$ is square-free
and $v_2(N)\leq8$. We let $R_T(D,M)$ denote the complete set of all norm-one equivalence
classes with characteristic matrix $T$, for the order of
discriminant $D$ and level $M$. When the order is clear we write $R_T$ for that set. If $R_T$ is nonempty, then
\begin{equation}\label{eq:n-flat}
 N^\flat\mid b^2-de.
\end{equation}
Moreover, uniformly for $d,e\asymp\Delta$, $|b|\ll_{r_0}\Delta$, and
$\Delta\leq N^A$,
\begin{equation}\label{eq:rt-bound}
 |R_T|\ll_{A,r_0,\eps}N^\eps.
\end{equation}
\end{lemma}

\begin{proof}
Write $M=2^kM_0$ with $M_0$ odd.  For $k>0$, let $R_{DM_0}$ be the Eichler
overorder of $R_N$ of level $M_0$; admissibility gives $2\nmid D$.  Then
\[
 L_N\subset L_{DM_0},\qquad \Gamma_N\subset\Gamma_{DM_0}.
\]
Thus every $\Gamma_{DM_0}$-orbit splits into at most
$[\Gamma_{DM_0}:\Gamma_N]$ many orbits of $\Gamma_N$.  The local Eichler
index, equivalently the quotient of the exact volume formulas in
\cite[Equation (2.7)]{BruinierLiMoller}, is
\begin{equation}\label{eq:two-index}
 |R_T(D,M)|\leq[\Gamma_{DM_0}:\Gamma_N]|R_T(D,M_0)|,
 \qquad
 [\Gamma_{DM_0}:\Gamma_N]=3\cdot2^{k-1}\leq384.
\end{equation}
It remains to treat the odd-level order.  Put
$N_0=DM_0$, $R=R_{N_0}$, $L=L(R)$, and
$\Gamma=R^1/\{\pm1\}$ for the rest of this proof.  Notice that the odd
part of $N_0$ is $N^\flat$.

Put $G=\operatorname{SO}(B^0)\simeq B^\times/\Q^\times$.  Fix a rational
pair $X_0=(x_0,y_0)$ with Gram matrix $T$, and, for each finite prime $p$,
write
\[
 \begin{aligned}
  \Kc_p&=\operatorname{im}(R_p^\times\longrightarrow G(\Q_p)),\\
  S_p(T)&=\{g\Kc_p:g^{-1}X_0\in L_p^2\},\\
  c_p(T)&=\#S_p(T).
 \end{aligned}
\]
The action of $G(F)$ on ordered pairs of Gram matrix $T$ is simply
transitive for $F=\Q$ and $F=\Q_p$.  Indeed, an isometry between the two
nondegenerate planes extends by Witt's theorem, its determinant can be
corrected on the orthogonal line, and an element of $G(F)$ fixing the pair
also fixes that line because it has determinant one.  The orbit map is a
homeomorphism: its inverse is obtained by completing $(x,y)$ with their
cross-product vector, and its only denominator is a fixed nonzero multiple
of $\det T$.  Hence the inverse image of the closed integral Gram fiber in
the compact set $L_p^2$ is compact.  Since $\Kc_p$ is open, every
$S_p(T)$ is finite.

We next record the global-to-local injection.  If $\Kc=\prod_p\Kc_p$, then
\begin{equation}\label{eq:global-unit-intersection}
 G(\Q)\cap\Kc=\operatorname{im}(R^\times).
\end{equation}
To see this, represent a class in the intersection by $q\in B^\times(\Q)$
and write $q=\lambda_pu_p$ locally, with
$\lambda_p\in\Q_p^\times$ and $u_p\in R_p^\times$.  Since
$v_p(\nrd q)=2v_p(\lambda_p)$, multiplication by the rational number having
valuation $-v_p(\lambda_p)$ at every $p$ gives an element, and its inverse,
in every $R_p$.  The identity
$R=B(\Q)\cap\prod_pR_p$ proves \eqref{eq:global-unit-intersection}.
Simple transitivity now maps full-unit global orbits injectively into the
restricted product $\prod_pS_p(T)$.  Passing from full units to norm-one units costs an index
at most two, and hence
\begin{equation}\label{eq:local-coset-product}
 |R_T(D,M_0)|\leq2\prod_pc_p(T).
\end{equation}

We count the local cosets directly.  At an odd prime, $L_p=R_p^0$.
If $B_p$ is split and $R_p$ is maximal, then $G(\Q_p)/\Kc_p$ is the
vertex set of the Bruhat--Tits tree.  An integral change in the two pair
coordinates puts the first elementary divisor first.  If it is a unit,
choose the first vector $x$ with $Q(x)=\epsilon\in\Z_p^\times$.  When
$\epsilon$ is nonsquare, $\Z_p[x]$ is the maximal unramified quadratic
order and its stable lattices give at most one homothety class.  When
$\epsilon=s^2$, the integral idempotents put
\[
 x=\begin{pmatrix}s&0\\0&-s\end{pmatrix},\qquad
 y=\begin{pmatrix}0&r\\u&0\end{pmatrix}.
\]
The $x$-stable lattices are
$\Lambda_n=\Z_pe_1\oplus p^n\Z_pe_2$, and $y$ stabilizes $\Lambda_n$
exactly for
\[
 -v_p(r)\leq n\leq v_p(u).
\]
As $v_p(ru)=t:=v_p(\det T)$, this gives at most $t+1$ vertices.
If the first elementary divisor is $p$, square-freeness of $d,e$ allows us
to choose $Q(x)=p\epsilon$.  The order $\Z_p[x]$ is then the maximal order
of a ramified quadratic field.  Its fractional ideals, modulo central
homothety, have two parity classes, so $c_p(T)\leq2$; here $t\geq2$.

If $p$ is an odd division prime, the unique maximal order and its Gross
lattice are invariant under conjugation, while
\[
 B_p^\times/(\Q_p^\times R_p^\times)\simeq\Z/2\Z.
\]
Thus a nonempty support has $c_p(T)=2$.  Reduction of $R_p^0$ modulo
the radical is one-dimensional, so such a represented pair has
$p\mid\det T$.  Finally, when $p\mid M_0$, the group $\Kc_p$ fixes an
oriented edge.  Since $(de,N)=1$, the maximal-order support is either one
vertex or an interval of $t+1$ vertices as above.  The former gives no
edge, and the latter gives $2t$ oriented edges.  Hence
\begin{equation}\label{eq:odd-local-counts}
 c_p(T)\leq
 \begin{cases}
  t+1,&p\nmid DM_0,\text{ unit first divisor},\\
  2,&p\nmid DM_0,\text{ first divisor }p,\\
  2,&p\mid D,\\
  2t,&p\mid M_0.
 \end{cases}
\end{equation}
These tree and order identifications are standard; see
\cite[Equation (23.2.4), Theorem 13.3.11, Lemma 23.5.15, and Sections
23.5.11--23.5.16]{Voight}.

It remains to count at $2$.  Since $-d,-e$ are odd fundamental
discriminants and $b$ is odd, $d\equiv e\equiv3\pmod4$ and
$t=v_2(\det T)\geq2$.  An integral shear gives
\[
 T\sim\operatorname{diag}(\epsilon_1,2^t\epsilon_2),
 \qquad \epsilon_1\equiv1\ \hbox{or}\ 5\pmod8.
\]
In the split algebra the Gross lattice of the standard maximal order is
\[
 L_2=\left\{\begin{pmatrix}z&2v\\2u&-z\end{pmatrix}:
 z,u,v\in\Z_2\right\},\qquad Q=z^2+4uv.
\]
For a unit-norm first vector $x$, membership in the Gross lattice of a
maximal order $O'$ is equivalent to $(1+x)/2\in O'$.  Indeed, writing
$x=2a-\trd(a)$, the oddness of $Q(x)=\trd(a)^2-4\nrd(a)$ forces
$\trd(a)$ to be odd; the converse follows from trace one.  If
$\epsilon_1\equiv5\pmod8$, this element generates the unramified maximal
quadratic order, whose stable lattices form one homothety class.  If
$\epsilon_1\equiv1\pmod8$, scale by a unit and take
$x=\operatorname{diag}(1,-1)$.  The integral idempotent $(1+x)/2$ forces
every stable lattice to split along the two eigenlines.  Writing
$y=\left(\begin{smallmatrix}0&r\\u&0\end{smallmatrix}\right)$, the initial
Gross integrality gives $r,u\in2\Z_2$, while
$Q(y)=ru=2^t\epsilon_2$.  Gross integrality at $\Lambda_n$ is equivalent to
\[
 1-v_2(r)\leq n\leq v_2(u)-1,
\]
which gives at most $t-1$ vertices.  In the dyadic division algebra the
unique maximal order is again invariant under conjugation and the
projective unit quotient has order two.  Consequently
\begin{equation}\label{eq:dyadic-local-counts}
 \begin{aligned}
  c_2(T)&\leq1 &&(\epsilon_1\equiv5\bmod8),\\
  c_2(T)&\leq t-1 &&(\epsilon_1\equiv1\bmod8),\\
  c_2(T)&=2 &&(2\mid D).
 \end{aligned}
\end{equation}

Every nontrivial local factor is therefore attached to a prime divisor of
$\det T$ and, in every branch of
\eqref{eq:odd-local-counts}--\eqref{eq:dyadic-local-counts},
\[
 c_p(T)\leq\bigl(v_p(\det T)+1\bigr)^2.
\]
At an odd prime dividing $D M_0$, the same rank-one reductions used above
also give $p\mid\det T$, proving $N^\flat\mid b^2-de$.  Thus
\[
 |R_T(D,M_0)|\leq2\prod_{p\mid\det T}
 \bigl(v_p(\det T)+1\bigr)^2=2\dd(|\det T|)^2.
\]
Together with \eqref{eq:two-index}, the bounds
$|\det T|\ll_{r_0}\Delta^2$ and $\Delta\leq N^A$, and the standard divisor
bound, this proves \eqref{eq:rt-bound}, uniformly in the admissible
factorization.
\end{proof}

\begin{proof}[Proof of \cref{thm:pair-intro}] For each ordered pair of quotient centers, we will need to choose lifts that realize the
respective quotient distance after shifting one of the two by a member of
$\Gamma_N$. In turn, we see that the corresponding simultaneous $\Gamma_N$-class of
primitive vectors is counted by the appropriate $R_T$. By \cref{lem:normalization}, we see that its characteristic matrix is
\[
 T=\begin{pmatrix}-d&b\\ b&-e\end{pmatrix}.
\]
We need to select signs for the two primitive vectors independently in such a way
that $b>0$.
Then, from the hyperboloid calculation preceding \cref{lem:normalization}, we can
read off $\cosh d(P,Q)=b/\sqrt{de}$.  Also, we can switch either sign, replacing $b$ by $-b$, but the CM center
remains unchanged. 
Then we will find that the sign-invariant form is
\begin{equation}\label{eq:distance-b} \cosh d(P,Q)=\frac{|b|}{\sqrt{de}}.
\end{equation}
We let $R_T$ denote the complete set of norm-one equivalence classes with
characteristic matrix $T$, and we write $N^\flat=N/2^{v_2(N)}$. We will treat the
singular and the nonsingular cases separately.

If $T$ is singular, then we can see that \eqref{eq:distance-b} gives distance zero. Since $d,e$ are both fundamental, we see that $b^2=de$ also gives $d=e$. Hence
we can conclude that the corresponding quotient centers agree. With $m(P)$ the multiplicity of a center, we can read off from the bounded
multiplicity that $\sum_Pm(P)^2\leq B_0\sum_Pm(P)=B_0|\mathscr P_\Delta|$. Together with the bounded fibres in \cref{rem:pair-formula}, we see that this
bound provides a contribution given by $O_{B_0}(|\mathscr P_\Delta|)$. 
For a nonsingular pair we will need \cref{lem:local-factors}. If the pair is at distance at most $r_0$, then we see that
\eqref{eq:distance-b} gives $O_{r_0}(\Delta)$ possibilities for $b$. We can now put
\[
 m=\frac{b^2-de}{N^\flat}.
\]
The form $Q=-\nrd$ has signature $(2,1)$. The reverse
Cauchy--Schwarz inequality on its negative cone gives
\[
 B_Q(x,y)^2\geq4Q(x)Q(y)=4de,
\]
with equality only when $x$ and $y$ are proportional. Since
$b=B_Q(x,y)/2$ and $T$ is nonsingular, the inequality is strict here;
therefore $b^2-de>0$. Together with \eqref{eq:n-flat}, this shows that
$m$ is a positive integer. Since the standing hypothesis gives
$N/N^\flat\leq2^8$, Equations
\eqref{eq:distance-b} and \eqref{eq:n-flat} give
\[
 1\leq m\ll_{r_0}1+\frac{\Delta^2}{N}.
\]
For fixed $b$ and $m$, we see that the product $de=b^2-N^\flat m$ is fixed as well. Then we can bound the number of ordered factor pairs $d,e\in[\Delta,2\Delta]$ by
its divisor function value. That divisor function value is $N^\eps$ throughout the
polynomial range. The conditions on fundamental discriminants and coprimality thin this set of
pairs, but as we can see they add nothing. The summation of \eqref{eq:rt-bound} over $b$ and $m$ therefore leads us to
\[ \ll_{A,r_0,\eps} \Delta\left(1+\frac{\Delta^2}{N}\right)N^\eps \] off the
diagonal. Restricting to submultisets removes pairs and does nothing else, apart from
the bounded multiplicities that were already controlled. This
proves the theorem.
\end{proof}

\section{A first absolute CM-packet moment}\label{sec:first-moment}

We now place \cref{thm:pair-intro} in Pasten's family. For an abc triple,
let $N$ be the conductor of the Frey curve attached to that triple, and
choose an admissible factorization $N=DM$. Let $0<\delta<1/10$ be a
parameter to be fixed at the end.  For this Frey curve, Pasten records
\[
 N=2^r\prod_{\substack{p\mid abc\\p>2}}p,
 \qquad 0\leq r\leq8;
\]
at every odd prime $p\mid abc$ one has $v_p(N)=1$
\cite[Section 3 and proof of Theorem 16.8]{Pasten}.  Thus the odd part of
$N$ is square-free and the hypotheses of \cref{sec:setup} hold for every
admissible factorization used below.  Pasten's sieve \cite[Theorem 11.1 and
Corollary 11.2]{Pasten}, along with
\begin{equation}\label{eq:parameters}
 x=N^{2/3+\delta},\qquad \vartheta=\delta/5,
\end{equation}
produces $H$ odd fundamental discriminants $-d$, with $x<d\leq2x$, in which
all the primes dividing $D$ are inert while all the primes dividing $M$ are
split in $\Q(\sqrt{-d})$, and
\begin{equation}\label{eq:sieve}
 H\gg x^{1-\delta},\qquad
 d^{1/2-\vartheta}<h_d\ll_\eps x^{1/2+\eps},\qquad
 \frac{L'}L(1,\chi_d)\ll\log\log d.
\end{equation}
In particular, all these discriminants are coprime to $N$.  In this situation we can apply the corollary that we cited, because
$x=N^{2/3+\delta}>N^{1/2+\delta}$, and we can note that
its constants are uniform in the admissible factorization $N=DM$.  Let $\mathscr C_d$ be the
full $2h_d$-element Galois multiset in
\cref{rem:galois-packet}, and set
\[
 \mathscr P=\bigsqcup_d\mathscr C_d,\qquad P=|\mathscr P|.
\]
Then $P\ll Hx^{1/2+\eps}$.  Put
$C(\mathscr P):=\#\{(P,Q)\in\mathscr P^2:d(P,Q)\leq2r_1\}$, where $r_1$ is
the fixed radius in \cref{lem:local-l1}.  Theorem~\ref{thm:pair-intro} gives
\begin{equation}\label{eq:pair-pasten}
 C(\mathscr P)\ll N^\eps\left(P+x+\frac{x^3}{N}\right).
\end{equation} \begin{lemma}[Uniform local $L^1$
evaluation]\label{lem:local-l1} Let $f(z)\,dz$ be a holomorphic differential
on $X_N$, and we put
\[
 F_N(z)=\frac{\sqrt N}{\|f\|_{\mathrm{Pas}}}\,y f(z),
 \qquad
 \|f\|_{\mathrm{Pas}}^2=
 \int_{X_N}|yf(z)|^2\,d\mu_{\mathrm{phys}}(z).
\]
We regard $F_N$ as a unitary automorphic section.  Its pointwise norm $|F_N|$ is a well defined function on $X_N$, and we note
that every evaluation and every integral below involving $F_N$ refers either
to this norm or to its powers.
For a hyperbolic radius $r_1$ that is fixed and sufficiently small, we have
\[
 |F_N(P)|\ll V_N\int_{B(P,r_1)}|F_N(z)|\,d\mu_{\mathrm{prob}}(z),
\]
uniformly for the levels in \cref{sec:setup}.  Moreover
we have $\|F_N\|_{L^2(\mu_{\mathrm{prob}})}^2=N/V_N$.
\end{lemma}

\begin{proof}
Fix a lift $\widetilde P=x_0+iy_0$.  Using the holomorphic submean inequality on a Euclidean disk of radius $\rho y_0$,
we read off, on a comparable hyperbolic disk, \[ |y_0f(\widetilde P)|\ll
\int_{B(\widetilde P,r_1)}|yf(z)|\,d\mu_{\mathrm{phys}}(z), \] since
$y\asymp y_0$ and $d\mu_{\mathrm{phys}}=dx\,dy/y^2$ hold there.

Next we need to make the descent to the orbifold uniform. The groups $\Gamma_N$ are cocompact, so we conclude that they contain no
parabolic elements. Also, a nonelliptic element has an integral reduced trace whose absolute
value is at least $3$, and so we conclude that its translation length is at
least $\ell_0=2\operatorname{arcosh}(3/2)$. We let $\mu_2$ denote a Margulis constant for the hyperbolic plane, and
choose $r_1<\min\{\mu_2/2,\ell_0/4\}$. Applying the two-dimensional Margulis lemma \cite{Ratcliffe} and then the
classification of elementary Fuchsian groups, we conclude that the subgroup
generated by \[ \{\gamma\in\Gamma_N: d(\widetilde P,\gamma\widetilde P)<2r_1\}
\] is finite. In fact, we note that an infinite elementary subgroup would either contain a
parabolic element or preserve a geodesic. In the latter case, for the subgroup generated by the displayed set to have
infinitely many elements, we need either a hyperbolic generator or two
distinct half-turn generators. As for the hyperbolic generator, it would contradict the
length gap, whereas we can also observe here that the product of the two
distinct half-turn generators has a translation length shorter than
$4r_1<\ell_0$. 
Thus the small-displacement subgroup is finite. A finite subgroup of
$\operatorname{PSL}_2(\R)$ is cyclic. Also, any nontrivial torsion lift in this situation has its
integral reduced trace lying in $\{-1,0,1\}$, so its projective order is no greater than $3$. As a consequence, we can say that the projection of $B(\widetilde P,r_1)$ to
$X_N$ has a uniformly bounded multiplicity, because an overlap of $B(\widetilde P,r_1)$ with its $\gamma$-translate forces
the inequality $d(\widetilde P,\gamma\widetilde P)<2r_1$.
Thus, when we descend the submean inequality that precedes this one, we can
see that the effect is that of changing just an absolute constant. Finally, we can take $d\mu_{\mathrm{phys}}=V_Nd\mu_{\mathrm{prob}}$ together
with the definition of $F_N$, and in this way we read off the
estimate that was asserted. The $L^2$ identity is immediate from the
normalizations.
\end{proof}

Define
\[
 T_d(F_N)=\frac1{2h_d}\sum_{P\in\mathscr C_d}|F_N(P)|
\]
and the sampling multiplicity
$\mu(z)=\sum_{P\in\mathscr P}\mathbf 1_{B(P,r_1)}(z)$.  Since any two intersecting balls have their centers at a bounded distance from one another, we can see that
\[
 \int\mu^2\,d\mu_{\mathrm{prob}}\ll \frac{C(\mathscr P)}{V_N}.
\] Indeed, after we have expanded the square, we will find that the pairs of
centers that contribute are just those whose distance apart is at most
$2r_1$, and we can also see that each ball intersection has a probability
measure that is equal to $O(1/V_N)$.
If we now sum \cref{lem:local-l1}, and then make use of the lower bound for $h_d$, and after that apply Cauchy--Schwarz in space, what we read off is the exact volume cancellation, and it comes out in the form
\begin{align}
 \frac1H\sum_dT_d(F_N)
 &\ll\frac{V_N}{Hx^{1/2-\vartheta}}
 \int\mu|F_N|\,d\mu_{\mathrm{prob}}\notag\\
 &\ll\frac{\sqrt{NC(\mathscr P)}}{Hx^{1/2-\vartheta}}.
 \label{eq:first-moment-basic}
\end{align}

If we now insert \eqref{eq:pair-pasten} together with the bound $P\ll Hx^{1/2+\eps}$, what we are left with is three terms:
\[
 E_{\mathrm{diag}}=N^{1/2}H^{-1/2}x^{-1/4+\vartheta+\eps},\quad
 E_x=N^{1/2}H^{-1}x^{\vartheta+\eps},\quad
 E_{\mathrm{off}}=H^{-1}x^{1+\vartheta+\eps}.
\]
Under \eqref{eq:parameters}--\eqref{eq:sieve}, we can compute that the $N$-exponents of these three are
\[
 -\frac{17}{60}\delta+\frac7{10}\delta^2,\qquad
 -\frac16-\frac15\delta+\frac65\delta^2,\qquad
 \tau=\left(\frac23+\delta\right)\frac{6\delta}{5}.
\]
It is the last of these three terms that is the limiting one here, and in this way we obtain the following proposition.

\begin{proposition}\label{prop:first-moment}
In the notation that we have set up in the preceding discussion, uniformly
in the admissible factorization $N=DM$,
\[
 \frac1H\sum_dT_d(F_N)\ll N^{\tau+\eps},
 \qquad
 \tau=\left(\frac23+\delta\right)\frac{6\delta}{5}.
\]
\end{proposition}

\newpage
\section{The arithmetic conclusion}\label{sec:pasten}

First we will need to choose a suitable packet on which logarithms can be taken.  Using \cite[Lemma 14.3]{Pasten}, we can see that if the corresponding rational
differential vanishes at one of the CM conjugate points, then it vanishes on
the entire Galois orbit as well.  Also, we can note that, by the same lemma, the associated modular form has at
most $N^{1+\eta}$
zeros on $\Hh$, under $\Gamma_N$-equivalence, for each fixed $\eta>0$.  Thus, for $\eta=\delta/5$, we can estimate the number $B$ of bad discriminants
as
\[
 B\ll N^{1+\delta/5}x^{-1/2+\delta/5},
 \qquad
 \frac BH\ll N^{-\delta/2+6\delta^2/5}=o(1).
\]
Since the negative line of a CM center determines the stabilizing torus of
that CM center, which also determines the imaginary quadratic field, we see
that all of the CM orbits corresponding to different quadratic fields are pairwise
disjoint. Using Markov's inequality together with the estimate from
\cref{prop:first-moment}, we can get a positive proportion of quadratic fields
with $T_d(F_N)\ll N^{\tau+\eps}$. Hence we can be sure that one of them is also zero-free. We now fix
such a field $K$. 
In fact, this field belongs to the original sieve, and we can also note that it
keeps the logarithmic-derivative bound that was recorded in \eqref{eq:sieve}.

In what follows we will let $f(z)\,dz=f_{D,M}(z)\,dz$ be the integral
quaternionic differential in Pasten's construction, that is to say, the
differential which comes from the modular parametrization in the form in which
he records it in his Equation (16.3). 
For all of the conjugates $P_\sigma$, and with a lift $\tau_\sigma\in\Hh$ chosen
for each one of them, we will write
$u_\sigma=|f(\tau_\sigma)|\operatorname{Im}\tau_\sigma>0$.
Then, from the normalization that was recorded in \cref{lem:local-l1}, we can
read off $u_\sigma=\|f\|_{\mathrm{Pas}}|F_N(P_\sigma)|/\sqrt N$.
Now, using ordinary AM--GM together with the first-moment estimate that was
selected above, we can then go on to obtain
\begin{equation}\label{eq:amgm}
 \frac1{2h_K}\sum_\sigma\log u_\sigma
 \leq \log\|f\|_{\mathrm{Pas}}-\frac12\log N
 +(\tau+\eps)\log N+O(1).
\end{equation}

Also, Proposition 14.2 in Pasten is given with respect to exactly this
hyperbolic metric, and therefore we can see that it implies
\[
 h_{\mathrm{Ar}}(P_K)\geq-\log(2M)
 -\frac1{2h_K}\sum_\sigma\log u_\sigma.
\]
Now, using his height formula \cite[Theorem 13.2]{Pasten} in the form in which
it is recorded there, after the functional equation has been applied, we can go
on to obtain
\[
 h_{\mathrm{Ar}}(P_K)
 =\frac{L'}L(1,\chi_K)+\frac12\log(dD)+O(1).
\]
Then, combining these formulas, and using $N=DM$ together with
$d\leq2N^{2/3+\delta}$, and together with the logarithmic-derivative estimate
that was retained above, we can now go on to deduce that
\begin{equation}\label{eq:petersson}
 -\log\|f\|_{\mathrm{Pas}}
 \leq\frac12\log M+
 \left(\frac13+\frac\delta2+\tau+\eps\right)\log N.
\end{equation}

In preparation for the final step of the argument, we will now single out the
one consequence of Pasten's modular-degree comparison that is used there.

\begin{proposition}[Pasten interface]\label{prop:pasten-interface}
For the Frey branch and for an admissible factorization $N=DM$ which has the
property that $M$ is divisible by at least two odd primes, we have the
following estimate
\[
 \prod_{p\mid D}v_p(\Delta_E)
 \ll_\eps N^{1+\eps}\|f\|_{\mathrm{Pas}}^{-2},
\]
uniformly in the admissible factorization.
\end{proposition}

\begin{proof}
Now, from Pasten's Theorem 6.1(b.2), what we can read off is
\[
 \log\prod_{p\mid D}v_p(\Delta_E)
 \leq\log\delta_{1,N}-\log\delta_{D,M}+O(\omega(D)).
\]
Using the degree relations (16.1)--(16.3), which are combined in
\cite[Theorem 16.1, pp.~305--307]{Pasten} and invoked in his Theorem
16.5(ii), we can then go on to read off the estimate
\begin{align*}
 \log\prod_{p\mid D}v_p(\Delta_E)
 &\leq2\log|c_f|+2\log\|f_{1,N}\|_2-2\log\|f_{D,M}\|_2\\
 &\quad+O(\log\log N+\omega(D)).
\end{align*}
Now the minimal isogeny has degree at most $163$ by an
application of the Mazur--Kenku theorem, in the form in which it is recorded
in \cite[proof of Theorem 16.1]{Pasten}.  Thus the difference between the two terms which are
associated with the Faltings height is $O(1)$.  In the same normalization, the classical estimate of
$\|f_{1,N}\|_2^2\ll N\log N$, the Manin constant estimate that is due to
Pasten \cite[Theorem 1.1 and proof of Theorem 16.1]{Pasten}, and the bound
$\omega(D)\ll\log N/\log\log N$ together imply that for sufficiently large
$N$ the right-hand side is at most
\[
 (1+\eps)\log N-2\log\|f_{D,M}\|_2.
\]
Finally, the non-normalized Petersson norms of Pasten take
into account the same physical hyperbolic measure \cite[Section
4.6]{Pasten}, and therefore we can conclude that
$\|f_{D,M}\|_2=\|f\|_{\mathrm{Pas}}$. All constants just used are absolute
or depend only on $\eps$ (the Manin-constant theorem is applied with the
fixed exceptional set $\{2\}$), so the estimate is uniform in $(D,M)$.
This proves the statement while additionally showing explicitly that the
Frey branch contributes no further factor of $D$.
\end{proof}

For every factorization satisfying the hypothesis of
\cref{prop:pasten-interface}, applying it to \eqref{eq:petersson} gives us the
cell estimate
\begin{equation}\label{eq:cell}
 \prod_{p\mid D}v_p(\Delta_E)
 \ll_\eps N^{q_\delta+\eps}M,\qquad
q_\delta=1+\left(\frac23+\delta\right)+2\tau
 =\frac53+\frac{13}{5}\delta+\frac{12}{5}\delta^2.
\end{equation}
The implied constant in \eqref{eq:cell} is uniform in the admissible
factorization.

It remains to distribute the factor $M$. In what follows we write the odd part of $N$ as $p_1\cdots p_m$, but the
bounded factor $2^r$ remains. If $m$ exceeds a fixed cutoff, then we can be sure that all subsequent
levels contain at least two odd primes, as required in
\cref{prop:pasten-interface}. In the even case
$m-2$ is even, and in the odd case $m-3$ is even. As a result, each $D_j$ described below is an admissible indefinite
quaternion discriminant and after the fixed cutoff has been increased, will
exceed 1. Moreover $(D_j,M_j)=1$, and all
constants are uniform in $j$. 
We then take each entry $(D_j,M_j)$ and apply to it the sieve,
the pair estimates, the first-moment selection, and the Pasten interface. Note that the selected field can depend on $j$ and that the earlier
established uniformity makes the cell estimate independent of $j$. If $m$ is even, apply \eqref{eq:cell}
cyclically with
\[
 M_j=2^rp_jp_{j+1},\qquad
 D_j=\frac{p_1\cdots p_m}{p_jp_{j+1}}.
\]
We can note that all valuations occur on $m-2$ left-hand sides, and that all
odd primes are included in two $M_j$. If $Q_{\mathrm{odd}}=\prod_{i=1}^m\bigl(2v_{p_i}(abc)\bigr)
=\prod_{i=1}^m v_{p_i}(\Delta_E)$, let $C_{\delta,\eps}\geq1$ be the uniform implied
constant in \eqref{eq:cell}.  Multiplication gives
\[
 Q_{\mathrm{odd}}^{m-2}
 \leq C_{\delta,\eps}^mN^{m(q_\delta+\eps)}\prod_{j=1}^mM_j,
 \qquad
 \prod_{j=1}^mM_j
 =2^{r(m-2)}N^2.
\]
Since $m/(m-2)\leq2$ for even $m\geq4$ and $r\leq8$, taking the
$(m-2)$-nd root leaves at most the bounded constant
$C_{\delta,\eps}^2 2^8$.  Thus
\begin{equation}\label{eq:even-prop}
 Q_{\mathrm{odd}}
 \ll N^{q_\delta+(2q_\delta+2)/(m-2)+O(\eps)}.
\end{equation}
If $m$ is odd, use cyclic triples
\[
 M_j=2^rp_jp_{j+1}p_{j+2},\qquad
 D_j=\frac{p_1\cdots p_m}{p_jp_{j+1}p_{j+2}},
\]
for which
\[
 Q_{\mathrm{odd}}^{m-3}
 \leq C_{\delta,\eps}^mN^{m(q_\delta+\eps)}\prod_{j=1}^mM_j,
 \qquad
 \prod_{j=1}^mM_j=2^{r(m-3)}N^3.
\]
Here $m/(m-3)\leq5/2$ for odd $m\geq5$, so the rooted constant is at
most $C_{\delta,\eps}^{5/2}2^8$.  Consequently
\begin{equation}\label{eq:odd-prop}
 Q_{\mathrm{odd}}
 \ll N^{q_\delta+(3q_\delta+3)/(m-3)+O(\eps)}.
\end{equation}
If $m\leq m_0$ for the chosen fixed cutoff, then we can see that
$\omega(abc)=m+1\leq m_0+1$. We can first employ Pasten's Proposition 15.1 with a smaller internal epsilon
and then select the maximum over these finitely many values, which gives
$\dd(abc)\ll_{\eps,m_0}R^{1+\eps}$ \cite[Proposition 15.1]{Pasten}. From his
Lemma 15.2 we can read off $v_2(abc)\ll_\eps R^\eps$ \cite[Lemma 15.2 and Equation
(16.4)]{Pasten}, and the Frey conductor satisfies $N\asymp R$. Finally,
\[
 P(abc)\leq\dd(abc)
 \leq\bigl(v_2(abc)+1\bigr)Q_{\mathrm{odd}}.
\]
In this last step we needed to use the inequality $v+1\leq2v$, and we needed
it for each positive integer value $v$ at the odd primes in question.

Given $\eps_0>0$, we first need to choose a fixed $0<\delta<1/10$ for which
$13\delta/5+12\delta^2/5<\eps_0/4$ holds. In the next step we need to ensure that the combined internal epsilon loss is below
$\eps_0/4$ and we need to select the prime-count cutoff such that both the overheads
in \eqref{eq:even-prop}--\eqref{eq:odd-prop} fall below $\eps_0/4$ above it. Here we can rely on Proposition 15.1 of Pasten, which takes care of those cases
where $m$ is equal to or lower than our cutoff. The finitely many conductors below the previously established thresholds end
up in the implied constant, and we do not track the effective dependence of
that constant on $\eps_0$. With that we have proved \cref{thm:main}.

\begin{remark}
We attain the exponent $5/3$ in the proof by means of the discriminant scale
$\Delta=N^{2/3+o(1)}$.  
We see that the above reasoning does not prove the abc conjecture.
In fact, we note that it handles only the product of valuation exponents and
the divisor function, not $abc$ itself.
\end{remark}


\section*{Notation}

\noindent{\setlength\tabcolsep{4pt}%
\begin{tabular}{@{}ll@{}}
$R=\rad(abc)$ & radical of the $abc$ triple\\
$P(abc)$, $\dd(n)$ & valuation product and divisor function\\
$B$ & indefinite quaternion division algebra over $\Q$\\ $B$ (in
\S\ref{sec:pasten}) & number of bad discriminants\\
$B^0$, $Q=-\nrd$ & trace-zero space and its quadratic form\\
$B_Q$ & full polarization of $Q$, so $B_Q(x,y)=\trd(xy)$\\
$D$ & discriminant of $B$\\
$M$ & Eichler level\\
$N=DM$ & reduced discriminant of $R_N$, and the Frey conductor\\
$R_N$ & Eichler order of reduced discriminant $N$\\
$\Gamma_N=R_N^1/\{\pm1\}$ & norm-one units modulo $\pm1$\\
$\Hh$ & upper half plane\\
$X_N=\Gamma_N\backslash\Hh$ & the Shimura curve\\
$V_N=\vol(X_N)$ & its hyperbolic volume\\
$L_N\subset B^0$ & Gross lattice of \cite[Equation
(2.9)]{BruinierLiMoller}\\
$G=\operatorname{SO}(B^0)$ & $\simeq B^\times/\Q^\times$\\
$\Kc_p$, $\Kc=\prod_p\Kc_p$ & local projective unit image and its product\\
$T$, $R_T(D,M)$ & characteristic matrix, and its norm-one class set\\
$S_p(T)$, $c_p(T)$ & local integral-coset set and its cardinality\\
$-d$, $-e$ & odd fundamental CM discriminants\\
$h_d$ & class number of $\Q(\sqrt{-d})$\\
$\mathscr C_d$ & the complete $2h_d$-element Galois multiset\\
$\Delta$ & the CM discriminant scale\\ $\mathscr P_\Delta$ & a multiset of
CM centers with $d\in[\Delta,2\Delta]$\\ $B_0$ & bound on multiplicities in
$\mathscr P_\Delta$\\
$\mathscr P=\bigsqcup_d\mathscr C_d$ & the full Galois CM multiset\\
$C(\mathscr P)$ & number of ordered pairs of $\mathscr P$ within $2r_1$\\
$H$ & number of discriminants produced by Pasten's sieve\\
$x=N^{2/3+\delta}$, $\vartheta=\delta/5$ & sieve parameters\\
$f(z)\,dz$ & holomorphic differential on $X_N$, $f=f_{D,M}$
in \S\ref{sec:pasten}\\
$\|f\|_{\mathrm{Pas}}$ & its non-normalized Petersson norm\\
$F_N$ & the unitary normalization $\sqrt N\,yf/\|f\|_{\mathrm{Pas}}$\\
$T_d(F_N)$ & first absolute moment of $|F_N|$ over $\mathscr C_d$\\
$\mu(z)$ & sampling multiplicity of the balls $B(P,r_1)$\\
$r_0$, $r_1$ & fixed hyperbolic radii\\
$\tau$, $q_\delta$ & moment exponent and cell exponent\\
$\Delta_E$ & discriminant of the Frey curve $E$\\
\end{tabular}}

\bibliographystyle{amsplain}
\bibliography{references}

\end{document}
